\documentclass[numsec,webpdf,modern,medium,namedate]{oup-authoring-template}

\onecolumn
\usepackage{amsmath,amssymb,amsthm,mathtools}
\usepackage{booktabs,graphicx}
\usepackage{microtype}
\graphicspath{{./}}

\theoremstyle{plain}
\newtheorem{theorem}{Theorem}
\newtheorem{proposition}{Proposition}
\newtheorem{lemma}{Lemma}
\newtheorem{corollary}{Corollary}
\theoremstyle{definition}
\newtheorem{assumption}{Assumption}
\newtheorem{definition}{Definition}
\theoremstyle{remark}
\newtheorem{remark}{Remark}
\newcommand{\ind}{\mathbf 1}

\begin{document}
\journaltitle{Journal of the Royal Statistical Society Series B: Statistical Methodology}
\DOI{}
\copyrightyear{2026}
\pubyear{2026}
\access{}
\appnotes{Original Article}
\firstpage{1}
\title[Inference for hard classification likelihoods]{Inference for Hard Classification Likelihoods: Moving-Boundary Curvature and Feasible Uncertainty Quantification}

\author[1,$\ast$]{Samyajoy Pal\ORCID{0000-0003-1339-0979}}
\author[1]{Rouven E. Haschka}
\authormark{Pal and Haschka}
\address[1]{\orgdiv{Chair of Data Analytics}, \orgname{Technical University of Kaiserslautern-Landau}, \orgaddress{\street{Gottlieb-Daimler-Stra\ss e 42}, \postcode{67663}, \state{Kaiserslautern}, \country{Germany}}}
\corresp[$\ast$]{Samyajoy Pal, Chair of Data Analytics, Technical University of Kaiserslautern-Landau, 67663 Kaiserslautern, Germany. \href{mailto:samyajoy.pal@rptu.de}{samyajoy.pal@rptu.de}}

\abstract{Hard classification replaces posterior membership probabilities by
deterministic assignments and can make heterogeneous mixture fitting markedly
simpler.  It also creates a nonsmooth criterion whose classification regions
move with the parameter.  Standard errors computed after freezing the fitted
labels ignore that motion.  We study classification maximum-likelihood
estimators as potentially misspecified M-estimators, with the population
classification functional---rather than the ordinary mixture parameter---as
their explicit target, and focus on a selected local optimum.
For regular decision hypersurfaces, we show that the
population Hessian contains a positive semidefinite surface integral in
addition to the usual within-class Hessian.  Because this term reduces the
positive curvature, fixed-label sandwich inference generally understates
uncertainty.  We establish consistency, a local stochastic expansion,
asymptotic normality, local-bootstrap validity, and consistency of a feasible
kernel estimator of the surface term.  An eigenvalue-constrained
stabilization is asymptotically inactive but prevents nonpositive estimated
curvature in finite samples.  Across 31,000 primary simulation attempts involving
heterogeneous Normal--Student-$t$ and Normal--skew-normal mixtures, the
correction materially improves standard-error calibration and confidence
interval coverage under overlap and imbalance, while becoming small under
separation.  Two Dry Bean analyses show both successful local calibration and
the distinct instability caused by optimization-basin selection; diagnostics
identify basin instability.  The method therefore supplies feasible local
inference together with diagnostics that identify when a local Gaussian
summary is inadequate.}
\keywords{classification likelihood; decision boundary; finite mixture; M-estimation; nonsmooth inference; sandwich covariance}
\maketitle

\section{Introduction}
\label{sec:introduction}

Many estimators combine smooth component criteria through a parameter-dependent
maximum.  Examples arise whenever observations are assigned to the currently
best-fitting group and the parameters are then estimated using those hard
assignments.  Such procedures are attractive computationally: conditional on
the assignments, a large optimization can separate into simpler component
problems.  Statistically, however, the resulting objective is only piecewise
smooth.  A small parameter perturbation moves the interfaces between the
regions, so differentiating the criterion as though its fitted partition were
fixed omits part of its curvature.
The paper's main message is that hard classification creates moving decision
boundaries, and these boundaries contribute first-order curvature even though
they have probability zero.  Mixtures provide the main implementation, but
the theory applies more generally to max-type classification likelihoods built
from smooth component criteria.

Classification expectation--maximization (CEM) provides an important instance
of this problem.  Finite mixtures are standard tools for density estimation,
clustering, and latent-population analysis \citep{mclachlan2000finite}.  The EM
algorithm maximizes the observed-data likelihood through probabilistic
memberships \citep{dempster1977maximum}; CEM instead inserts a classification
step and assigns every observation to its largest weighted component density
\citep{celeux1992classification}.  Once labels are fixed, component updates
are partly decoupled.  This is especially useful in heterogeneous mixtures,
where different components use different distribution families and some
updates require numerical optimization \citep{pal2026flexible}.

Model-based clustering has developed a broad theory of mixture fitting,
identifiability, parsimonious covariance modelling, initialization, and order
selection \citep{titterington1985statistical,redner1984mixture,
banfield1993model,fraley2002model,keribin2000consistent}. Mixture likelihoods
can be singular or multimodal, motivating constrained estimation and careful
optimization diagnostics \citep{hathaway1985constrained,chen1995optimal,
mclachlan2008em}. Partition assessment and selection add a complementary
literature based on adjusted agreement, silhouettes, gaps, entropy, and
resampling stability \citep{hubert1985comparing,rousseeuw1987silhouettes,
tibshirani2001gap,celeux1996entropy,benhur2002stability,fang2012selection}.
These concerns are relevant to our basin diagnostics, but are distinct from
estimating the local curvature of a prespecified classification-likelihood
target.

Classification likelihood predates CEM. Likelihood-ratio clustering was
studied by \citet{scott1971clustering}, and
\citet{bryant1978classification} derived early large-sample properties of
classification maximum-likelihood estimates. Later work used classification
criteria for mixture-based clustering and model selection, most notably the
integrated completed likelihood \citep{biernacki2000icl} and conditional
classification likelihood \citep{baudry2015conditional}. These contributions
clarify that classification is a statistical criterion in its own right, but
they do not provide a general feasible sandwich covariance that estimates the
curvature contributed by moving multivariate decision faces in heterogeneous
component models.

Comparisons of mixture and classification sampling schemes likewise show that
hard classification targets must not be interpreted automatically as
ordinary-mixture parameters \citep{ganesalingam1989classification}. Related
order-selection work shows that nontrivial classification uncertainty changes
the asymptotic behaviour of complete-likelihood criteria \citep{hui2015order}.
For distance-based hard clustering, early distributional and
significance-test analyses further emphasize the role of a well-separated
population optimum \citep{hartigan1978asymptotic,bock1985probabilistic}.

The computational simplification changes both the estimand and the inference.
The CEM output maximizes a classification likelihood, not the ordinary
observed-data mixture likelihood.  Its natural population target is therefore
the maximizer of the expected classification criterion, and need not equal the
parameter that generated an ordinary finite mixture.  Moreover, the common
practice of conditioning on the fitted labels treats a data-dependent
partition as fixed.  This discards the effect of observations that cross a
decision boundary under an arbitrarily small parameter perturbation.  The
effect is negligible only when the data-generating density is sufficiently
small on the active boundaries; it can be first-order under component overlap.

The closest geometric precedent is the asymptotic theory of $k$-means.
\citet{pollard1981strong} established consistency and
\citet{pollard1982clt} derived a central limit theorem in which the geometry of
moving Voronoi cells is fundamental. Our setting replaces squared-distance
cells by general weighted component log-density comparisons, permits different
families and misspecification, identifies the active-face contribution in a
coordinate-invariant surface form, and develops a sample-based estimator of
that contribution suitable for routine sandwich inference.

This phenomenon is easy to miss because a regular boundary has probability
zero.  Probability zero eliminates a boundary contribution to the first
derivative, but not to the second.  We show that differentiation of the
population maximum produces a surface integral over each active decision
hypersurface.  Its integrand combines the data density on the boundary, the
sensitivity of the pairwise score contrast to the parameter, and the spatial
gradient transverse to the boundary.  The term is positive semidefinite in the
population Hessian.  It consequently subtracts from the usual positive
fixed-class curvature and generally increases the sandwich variance.

Our first contribution is a theory of the classification maximum-likelihood
functional under possible model misspecification.  We establish existence and
consistency, derive the moving-boundary Hessian, and prove a uniform local
empirical-process expansion for the nonsmooth maximum criterion.  These
results yield root-$n$ consistency, asymptotic linearity, and normality for a
sufficiently accurate local maximizer.  We also state explicit conditions
under which an algorithmic stationary solution inherits that distribution and
under which a nonparametric bootstrap is valid.  This separation is essential:
asymptotic theory for a population target does not guarantee that a multistart
algorithm or every bootstrap refit reaches its basin.

Our second contribution is feasible inference without estimating the full
data-generating density or numerically integrating over a hypersurface.  By a
coarea identity, the surface curvature is equivalently a Dirac-delta moment of
the pairwise score contrasts.  We estimate it using a one-dimensional kernel
around each active contrast, obtaining a procedure whose observation-level
cost is linear in sample size and does not require a grid in the observation
space.  We derive its bias, variance, oracle limit distribution, plug-in rate,
and consistency.  Because sampling variability can make the corrected
curvature nonpositive in finite samples, we introduce an eigenvalue-constrained
stabilization and prove that it is asymptotically inactive.
The construction uses the coarea formula \citep{federer1969geometric} and
standard one-dimensional kernel principles \citep{silverman1986density,
wand1995kernel}; its plug-in analysis is framed as a stochastic-equicontinuity
problem rather than a full multivariate density-estimation problem
\citep{andrews1994empirical}.

Our third contribution is an empirical assessment designed around the actual
statistical target and the method's operational limitations.  The simulation
study uses independently approximated population classification targets rather
than the ordinary mixture-generating parameters, retains failed fits, varies
overlap, imbalance, dimension, number of components, and heterogeneous family,
and examines bandwidth sensitivity.  Across the main designs, fixed-label
inference is most anti-conservative when the boundary carries appreciable
mass; the proposed correction moves estimated standard errors and coverage
toward their empirical values.  Additional large-sample runs distinguish slow
convergence from persistent failure.  In two real-data analyses we compare the
analytic covariance with a local nonparametric bootstrap and separately use
unrestricted bootstrap refits to diagnose basin-selection instability.

The contribution differs from ordinary likelihood information calculations,
including missing-information decompositions such as \citet{Louis1982FindingTO}.
Here the issue is not integration over an unobserved label conditional on a
smooth likelihood.  It is the motion of deterministic classification regions
inside a max-type population criterion.  It also differs from applying a
generic sandwich formula after differentiation: the relevant derivative
matrix itself contains a singular surface contribution.  The results connect
classical misspecified M-estimation \citep{Huber1967,white1982maximum} with the
geometry of hard classification.

Technically, max-type criteria sit within the wider literature on nonsmooth
M-estimation and empirical processes. Depending on how sharply the population
criterion changes, nonsmooth estimators can retain a Gaussian root-$n$ limit
or exhibit cube-root behaviour \citep{kim1990cube,sherman1993limiting}; the
regular active-face and positive-curvature assumptions below identify the
former regime. Our bootstrap statement builds on classical nonparametric
bootstrap theory \citep{efron1979bootstrap,efron1993introduction}, while the
application diagnostics recognize that clustering instability is a separate
inferential object.

\subsection{A one-dimensional illustration of the boundary term}
\label{sec:one-dimensional-illustration}
Before introducing the general notation, a scalar example makes the source
and magnitude of the additional curvature transparent.  
Consider two components on the real line and suppose that their contrast
$g(x;\boldsymbol{\theta})=q_1(x;\boldsymbol{\theta})-
q_2(x;\boldsymbol{\theta})$ has one active regular boundary point $x_0$ at
$\boldsymbol{\theta}^*$.  The surface correction in
Theorem~\ref{thm:population-curvature} then reduces to
\begin{equation}
 \mathbf B=\frac{p_0(x_0)\nabla_{\boldsymbol\theta}g(x_0;\boldsymbol\theta^*)
 \nabla_{\boldsymbol\theta}g(x_0;\boldsymbol\theta^*)^\mathsf T}
 {|\nabla_x g(x_0;\boldsymbol\theta^*)|}.
 \label{eq:one-dimensional-boundary}
\end{equation}
Thus the correction is large when the data density is high at the decision
boundary, all else equal.  It is small when the boundary lies in a low-density
region, which explains why fixed-label inference can be adequate for strongly
separated groups.

\newpage
The general development below replaces the single point $x_0$ by active
hypersurfaces and turns this population identity into feasible inference.
Section~\ref{sec:methodology} defines the estimand, derives the active-boundary
curvature, establishes the limit theory, and constructs the feasible
covariance estimator. Section~\ref{sec:results} reports the simulation study
and real-data analyses. Sections~\ref{sec:discussion} and the Conclusion state
the implications and limitations. Technical proofs, complete empirical
summaries, direct Gaussian-limit diagnostics, failure accounting, and
application-level parameter tables are provided in the supplement.

\section{Methodology}
\label{sec:methodology}

\subsection{Classification likelihood and the population target}
\label{sec:framework}
Let $\mathbf{X}_1,\ldots,\mathbf{X}_n$ be independent copies of a random vector
$\mathbf{X}\in\mathcal X\subseteq\mathbb R^p$ with distribution $P_0$.  We do not
initially require $P_0$ to belong to a finite-mixture model.  This distinction
allows the classification likelihood to be treated as a possibly
misspecified criterion and makes its population target explicit.

Consider $k\geq2$ positive component scores
\begin{equation}
 q_j(\mathbf{x};\boldsymbol{\theta})=\log\pi_j(\boldsymbol{\theta})+\log f_j(\mathbf{x};\boldsymbol{\alpha}_j(\boldsymbol{\theta})),
 \qquad j=1,\ldots,k,
 \label{eq:component-score}
\end{equation}
where $f_j(\cdot;\boldsymbol{\alpha}_j)$ may belong to different parametric families.  The
component parameter $\boldsymbol{\alpha}_j$ has dimension $d_j$.  The full parameter
$\boldsymbol{\theta}$ is expressed in a local unconstrained coordinate system and takes
values in $\Theta\subset\mathbb R^d$, where
\begin{equation}
 d=(k-1)+\sum_{j=1}^k d_j.
 \label{eq:parameter-dimension}
\end{equation}
For example, the first $k-1$
coordinates may be multinomial-logit coordinates for the mixing weights,
while covariance matrices may be represented by unconstrained Cholesky
coordinates.  All derivatives below are with respect to this fixed coordinate
system.

Define
\begin{equation}
 m_{\boldsymbol{\theta}}(\mathbf{x})=\max_{1\leq j\leq k}q_j(\mathbf{x};\boldsymbol{\theta}),\qquad
 M_n(\boldsymbol{\theta})=\mathbb P_n m_{\boldsymbol{\theta}},\qquad
 M(\boldsymbol{\theta})=P_0m_{\boldsymbol{\theta}},
 \label{eq:cml-criteria}
\end{equation}
where $\mathbb P_n=n^{-1}\sum_{i=1}^n\delta_{\mathbf{X}_i}$.  The population
classification maximum-likelihood target is
\begin{equation}
 \boldsymbol{\theta}^*=\arg\max_{\boldsymbol{\theta}\in\Theta}M(\boldsymbol{\theta}).
 \label{eq:cml-target}
\end{equation}
The equality in \eqref{eq:cml-target} defines a statistical functional of
$P_0$; it does not assert that $\boldsymbol{\theta}^*$ is the parameter generating $P_0$.
Even under a correctly specified ordinary mixture model, the CML target and
the ordinary mixture-likelihood parameter need not coincide.

The theoretical estimator is any measurable approximate maximizer satisfying
\begin{equation}
 M_n(\widehat{\boldsymbol{\theta}}_n)
 \geq \sup_{\boldsymbol{\theta}\in\Theta}M_n(\boldsymbol{\theta})-r_n,
 \qquad r_n=o_p(1).
 \label{eq:approximate-maximizer}
\end{equation}
Exact empirical maximizers correspond to $r_n=0$.  Stronger control of $r_n$
will be imposed when deriving a first-order distribution.  Condition
\eqref{eq:approximate-maximizer} is a property to be verified for an
algorithmic output; calling an arbitrary stationary point a CML estimator does
not make the condition true.

For $j\ne l$, write
\begin{equation}
 g_{jl}(\mathbf{x};\boldsymbol{\theta})=q_j(\mathbf{x};\boldsymbol{\theta})-q_l(\mathbf{x};\boldsymbol{\theta}).
 \label{eq:pairwise-contrast}
\end{equation}
Let $a_{\boldsymbol{\theta}}(\mathbf{x})$ be the smallest index in
$\arg\max_jq_j(\mathbf{x};\boldsymbol{\theta})$.  This deterministic convention makes
$a_{\boldsymbol{\theta}}$ measurable and defines the disjoint classification cells
\begin{equation}
 \mathcal R_j(\boldsymbol{\theta})=\{\mathbf{x}:a_{\boldsymbol{\theta}}(\mathbf{x})=j\},\qquad j=1,\ldots,k.
 \label{eq:classification-cells}
\end{equation}
The convention has no effect on population quantities whenever ties have
$P_0$-probability zero.

\subsection{Parameter identifiability}

When two components use the same family, simultaneous permutations of their
weights and component parameters leave $m_{\boldsymbol{\theta}}$ unchanged.  We therefore
work on an identifiable representative set $\Theta$.  This can be obtained by
a scientifically meaningful component ordering, a fixed family assignment in
a heterogeneous model, or a canonical ordering of a component functional.
The chosen ordering must hold in a neighbourhood of $\boldsymbol{\theta}^*$; ordering
components after estimation is not a substitute for local identifiability.

The ordinary mixture density
$\sum_j\pi_jf_j(\cdot;\boldsymbol{\alpha}_j)$ is used for data generation in some
experiments, but it is not used to replace the unknown $P_0$ in the general
theory.  In particular, the feasible boundary estimator developed later is
based on observations from $P_0$ and does not assume that the mixture evaluated
at $\boldsymbol{\theta}^*$ equals the data-generating density.

\begin{assumption}[Basic criterion conditions]
\label{ass:basic-criterion}
The following conditions hold.
\begin{enumerate}
\item $\Theta$ is a compact, identifiable subset of $\mathbb R^d$.
 \item For every $j$, $(\mathbf{x},\boldsymbol{\theta})\mapsto q_j(\mathbf{x};\boldsymbol{\theta})$ is jointly
 measurable, and $\boldsymbol{\theta}\mapsto q_j(\mathbf{x};\boldsymbol{\theta})$ is continuous for
 $P_0$-almost every $\mathbf{x}$.
 \item There is a measurable function $F$ with $P_0F<\infty$ such that
 $\sup_{\boldsymbol{\theta}\in\Theta}\max_j|q_j(\mathbf{X};\boldsymbol{\theta})|\leq F(\mathbf{X})$ almost surely.
 \item $M$ has a unique, well-separated maximizer $\boldsymbol{\theta}^*$: for every
 $\varepsilon>0$,
 \[
  M(\boldsymbol{\theta}^*)-
  \sup_{\boldsymbol{\theta}\in\Theta:\,\|\boldsymbol{\theta}-\boldsymbol{\theta}^*\|\geq\varepsilon}M(\boldsymbol{\theta})>0.
 \]
\end{enumerate}
\end{assumption}

Compactness is used only for the global consistency argument.  Later local
results require instead that $\boldsymbol{\theta}^*$ be an interior point of an open
coordinate neighbourhood and that the empirical procedure enters that
neighbourhood with probability tending to one.


\subsection{Local geometry of active decision boundaries}
\label{sec:geometry}
The equality set $\{\mathbf{x}:g_{jl}(\mathbf{x};\boldsymbol{\theta})=0\}$ is not necessarily a boundary of
the classification partition.  A third component can dominate both $j$ and
$l$ on part of that set.  For $j<l$, define the relatively open active face
\begin{equation}
 \mathcal F_{jl}(\boldsymbol{\theta})
 =\left\{\mathbf{x}:q_j(\mathbf{x};\boldsymbol{\theta})=q_l(\mathbf{x};\boldsymbol{\theta})>
 q_m(\mathbf{x};\boldsymbol{\theta})\ \text{for every }m\notin\{j,l\}\right\}.
 \label{eq:active-face}
\end{equation}
Only $\mathcal F_{jl}(\boldsymbol{\theta})$ separates cells $j$ and $l$ locally.  Its
closure can meet other faces at multiple-tie junctions.  We write
$\mathcal F_{jl}^*=\mathcal F_{jl}(\boldsymbol{\theta}^*)$.

For a smooth scalar map $g$, the notation $\nabla_{\mathbf{x}}g$ and $\nabla_{\boldsymbol{\theta}}g$ denotes
gradients with respect to $\mathbf{x}$ and $\boldsymbol{\theta}$, respectively; $\nabla_{\boldsymbol{\theta}}^2q_j$ denotes
the Hessian with respect to $\boldsymbol{\theta}$.  Let $\mathcal H^{p-1}$ denote
$(p-1)$-dimensional Hausdorff measure.

\begin{assumption}[Regular active-boundary geometry]
\label{ass:boundary-geometry}
There is an open neighbourhood $U$ of the interior point $\boldsymbol{\theta}^*$ such that:
\begin{enumerate}
 \item $P_0$ has a Lebesgue density $p_0$ that is continuous on a
 neighbourhood of the closures of the active faces at $\boldsymbol{\theta}^*$.
 \item Each $q_j(\mathbf{x};\boldsymbol{\theta})$ is twice continuously differentiable in $\boldsymbol{\theta}$
 and continuously differentiable in $\mathbf{x}$ on the relevant neighbourhood.
 \item For every nonempty $\mathcal F_{jl}^*$,
 $\nabla_{\mathbf{x}}g_{jl}(\mathbf{x};\boldsymbol{\theta}^*)\ne0$ on its closure except possibly on a set of
 $\mathcal H^{p-1}$-measure zero.
 \item The set on which three or more component scores attain the common
 maximum has $\mathcal H^{p-1}$-measure zero.
 \item The derivatives can be passed through the volume and surface
 integrals below.  In particular, the interior Hessian terms are
 $P_0$-integrable and
 \[
  \int_{\mathcal F_{jl}^*}
  \frac{p_0(\mathbf{x})\|\nabla_{\boldsymbol{\theta}}g_{jl}(\mathbf{x};\boldsymbol{\theta}^*)\|^2}
       {\|\nabla_{\mathbf{x}}g_{jl}(\mathbf{x};\boldsymbol{\theta}^*)\|}
  \,d\mathcal H^{p-1}(\mathbf{x})<\infty
 \]
 for every $j<l$.
\end{enumerate}
\end{assumption}

Assumption~\ref{ass:boundary-geometry}(4) is a geometric condition, not merely
a zero-$P_0$-probability condition.  Every regular hypersurface has
$P_0$-probability zero under a Lebesgue density, while it can have positive
$\mathcal H^{p-1}$ measure and contribute to curvature.

At points with a unique winning component, define
\begin{equation}
 \boldsymbol{\psi}(\mathbf{x};\boldsymbol{\theta})
 =\sum_{j=1}^k
 \ind\{a_{\boldsymbol{\theta}}(\mathbf{x})=j\}\nabla_{\boldsymbol{\theta}}q_j(\mathbf{x};\boldsymbol{\theta}).
 \label{eq:classification-score}
\end{equation}
On ties the right-hand side uses the deterministic convention from
\eqref{eq:classification-cells}.  Its value on those points is immaterial under
Assumption~\ref{ass:boundary-geometry}.

\begin{theorem}[Population derivatives and active-boundary curvature]
\label{thm:population-curvature}
Suppose Assumption~\ref{ass:boundary-geometry} holds and the indicated
first-order derivatives admit integrable envelopes.  Then $M$ is differentiable
at $\boldsymbol{\theta}^*$ with
\begin{equation}
 \nabla_{\boldsymbol{\theta}}M(\boldsymbol{\theta}^*)=P_0\boldsymbol{\psi}(\mathbf{X};\boldsymbol{\theta}^*).
 \label{eq:population-gradient}
\end{equation}
It is twice differentiable at $\boldsymbol{\theta}^*$, and
\begin{align}
 \mathbf{H}:=\nabla_{\boldsymbol{\theta}}^2M(\boldsymbol{\theta}^*)
 &=P_0\left[
   \sum_{j=1}^k\ind\{a_{\boldsymbol{\theta}^*}(\mathbf{X})=j\}
   \nabla_{\boldsymbol{\theta}}^2q_j(\mathbf{X};\boldsymbol{\theta}^*)
  \right]+\mathbf{B},                                      \label{eq:hessian-decomposition}\\
 \mathbf{B}
 &=\sum_{1\leq j<l\leq k}
   \int_{\mathcal F_{jl}^*}
   \frac{p_0(\mathbf{x})\nabla_{\boldsymbol{\theta}}g_{jl}(\mathbf{x};\boldsymbol{\theta}^*)
        \nabla_{\boldsymbol{\theta}}g_{jl}(\mathbf{x};\boldsymbol{\theta}^*)^\mathsf T}
        {\|\nabla_{\mathbf{x}}g_{jl}(\mathbf{x};\boldsymbol{\theta}^*)\|}
   \,d\mathcal H^{p-1}(\mathbf{x}).                                \label{eq:boundary-curvature}
\end{align}
Consequently, if
\begin{equation}
 \mathbf{I}_c=-P_0\left[
   \sum_{j=1}^k\ind\{a_{\boldsymbol{\theta}^*}(\mathbf{X})=j\}
   \nabla_{\boldsymbol{\theta}}^2q_j(\mathbf{X};\boldsymbol{\theta}^*)
  \right],
 \label{eq:fixed-class-information}
\end{equation}
then the positive population curvature is
\begin{equation}
 \mathbf{A}=-\mathbf{H}=\mathbf{I}_c-\mathbf{B}.
 \label{eq:positive-curvature}
\end{equation}
In particular, $\mathbf{B}$ is positive semidefinite and parameter-dependent
classification weakens the positive curvature relative to fixed-label
analysis.
\end{theorem}

\begin{proof}
Away from ties, the maximizing index is locally constant in $\boldsymbol{\theta}$, so the
first derivative of $m_{\boldsymbol{\theta}}(\mathbf{x})$ equals the derivative of the winning score.
The envelope condition and the fact that regular ties have $P_0$-probability
zero give \eqref{eq:population-gradient} by dominated differentiation.

For the second derivative, first consider a neighbourhood containing only an
active face between $j$ and $l$.  Locally,
\[
 m_{\boldsymbol{\theta}}(\mathbf{x})=q_l(\mathbf{x};\boldsymbol{\theta})+\{g_{jl}(\mathbf{x};\boldsymbol{\theta})\}_+.
\]
In the distributional derivative with respect to $\boldsymbol{\theta}$, differentiating
the positive-part term a second time produces the ordinary within-region
Hessian and
\[
 \delta\{g_{jl}(\mathbf{x};\boldsymbol{\theta}^*)\}
 \nabla_{\boldsymbol{\theta}}g_{jl}(\mathbf{x};\boldsymbol{\theta}^*)\nabla_{\boldsymbol{\theta}}g_{jl}(\mathbf{x};\boldsymbol{\theta}^*)^\mathsf T.
\]
The coarea formula \citep{evans2015measure} converts its integral with respect to
$p_0(\mathbf{x})\,dx$ into the corresponding term in
\eqref{eq:boundary-curvature}.  A partition of unity over the unique-winner
regions and the relatively open active faces combines these local identities.
Inactive portions of pairwise equality sets do not enter because the maximum
is locally attained by another component there.  Multiple-tie junctions make
no contribution by Assumption~\ref{ass:boundary-geometry}(4).  Summing over
$j<l$ proves \eqref{eq:hessian-decomposition}; the remaining statements follow
from the definitions.
\end{proof}

\begin{remark}[Relation to $k$-means]
\label{rem:kmeans}
For $q_j(\mathbf{x};\boldsymbol{\theta})=-\|\mathbf{x}-
\boldsymbol{\mu}_j\|^2$, the active faces are Voronoi boundaries.  Pollard's
$k$-means boundary geometry is therefore recovered as a special
hard-classification case \citep{pollard1982clt}.  The result extends it to
weighted log-density comparisons, heterogeneous families, and misspecification.
\end{remark}

\begin{corollary}[Negligible-boundary limit]
\label{cor:negligible-boundary}
If $p_0=0$ $\mathcal H^{p-1}$-almost everywhere on every active face, then
$\mathbf{B}=0$ and fixed-classification curvature is correct to first order.
More generally, the size of the correction is governed by probability density
on the active faces, the sensitivity $\nabla_{\boldsymbol{\theta}}g_{jl}$, and the spatial
transversality $\|\nabla_{\mathbf{x}}g_{jl}\|$.
\end{corollary}


\subsection{Asymptotic theory}
\label{sec:theory}
We first establish consistency without imposing the smooth active-boundary
conditions needed for inference.

\begin{proposition}[Continuity and existence of the CML target]
\label{prop:population-existence}
Under Assumption~\ref{ass:basic-criterion}(1)--(3), $M$ is continuous on
$\Theta$ and attains its maximum.  Under part (4), the maximizer is the unique
target $\boldsymbol{\theta}^*$.
\end{proposition}

\begin{proof}
For $P_0$-almost every $\mathbf{x}$, the maximum of the finite collection of continuous
maps $\{q_j(\mathbf{x};\cdot):1\leq j\leq k\}$ is continuous.  The envelope in
Assumption~\ref{ass:basic-criterion}(3) and dominated convergence imply
continuity of $M$.  Existence follows from compactness of $\Theta$; uniqueness
is part (4).
\end{proof}

\begin{theorem}[Uniform convergence of the classification criterion]
\label{thm:uniform-convergence}
Under Assumption~\ref{ass:basic-criterion}(1)--(3),
\begin{equation}
 \sup_{\boldsymbol{\theta}\in\Theta}|M_n(\boldsymbol{\theta})-M(\boldsymbol{\theta})|\xrightarrow{p}0.
 \label{eq:uniform-convergence}
\end{equation}
If the observations form an infinite iid sequence, the convergence can be
taken almost surely.
\end{theorem}

\begin{proof}
The class $\{m_{\boldsymbol{\theta}}:\boldsymbol{\theta}\in\Theta\}$ has an integrable envelope.  It is
indexed by a compact finite-dimensional parameter set and is pointwise
continuous in $\boldsymbol{\theta}$ almost surely.  A finite covering argument, using
dominated continuity to control the oscillation within each cover element,
reduces the supremum to finitely many ordinary laws of large numbers.  This
gives \eqref{eq:uniform-convergence}; applying the strong law in the same
argument gives the almost-sure version.
\end{proof}

\begin{theorem}[Consistency of approximate CML maximizers]
\label{thm:cml-consistency}
Suppose Assumption~\ref{ass:basic-criterion} holds and
$\widehat{\boldsymbol{\theta}}_n$ satisfies \eqref{eq:approximate-maximizer}.  Then
\begin{equation}
 \widehat{\boldsymbol{\theta}}_n\xrightarrow{p}\boldsymbol{\theta}^*.
 \label{eq:cml-consistency}
\end{equation}
If the uniform convergence in Theorem~\ref{thm:uniform-convergence} and
$r_n\to0$ both hold almost surely, then the convergence in
\eqref{eq:cml-consistency} is almost sure.
\end{theorem}

\begin{proof}
Fix $\varepsilon>0$ and let
\[
 \Delta_\varepsilon=M(\boldsymbol{\theta}^*)-
 \sup_{\|\boldsymbol{\theta}-\boldsymbol{\theta}^*\|\geq\varepsilon}M(\boldsymbol{\theta})>0.
\]
On the event
\[
 2\sup_{\boldsymbol{\theta}\in\Theta}|M_n(\boldsymbol{\theta})-M(\boldsymbol{\theta})|+r_n
 <\Delta_\varepsilon,
\]
no approximate maximizer satisfying \eqref{eq:approximate-maximizer} can lie
outside the $\varepsilon$-ball around $\boldsymbol{\theta}^*$.  The probability of this
event tends to one by Theorem~\ref{thm:uniform-convergence} and $r_n=o_p(1)$.
The almost-sure statement follows identically.
\end{proof}

Theorem~\ref{thm:cml-consistency} deliberately separates statistical and
algorithmic claims.  A CEM iteration that terminates at an arbitrary local
maximum need not satisfy \eqref{eq:approximate-maximizer}.  Later we give a
local transfer result under consistent initialization and a quantified
optimization error.

For the first-order theory, Theorem~\ref{thm:population-curvature} supplies the
deterministic quadratic expansion
\begin{equation}
 M(\boldsymbol{\theta}^*+\mathbf{u})-M(\boldsymbol{\theta}^*)
 =\tfrac12\mathbf{u}^\mathsf THu+o(\|\mathbf{u}\|^2),
 \label{eq:population-quadratic-expansion}
\end{equation}
because $\nabla_{\boldsymbol{\theta}}M(\boldsymbol{\theta}^*)=0$ at an interior maximizer.  The next theoretical
layer will establish the matching local empirical-process expansion and show
that a sufficiently accurate local maximizer is asymptotically linear with
curvature $\mathbf{A}=\mathbf{I}_c-\mathbf{B}$.

\subsection{Local stochastic expansion}

Write $\mathbb G_n=\sqrt n(\mathbb P_n-P_0)$ and
$\boldsymbol{\psi}^*(\mathbf{x})=\boldsymbol{\psi}(\mathbf{x};\boldsymbol{\theta}^*)$.  Although $m_{\boldsymbol{\theta}}(\mathbf{x})$ is continuous in
$\boldsymbol{\theta}$, it is not differentiable when two component scores tie.  The
probability of observations lying in a shrinking tube around an active face
therefore controls the empirical remainder.

For $j<l$, define the activity margin
\begin{equation}
 \Delta_{jl}(\mathbf{x};\boldsymbol{\theta})
 =\min\{q_j(\mathbf{x};\boldsymbol{\theta}),q_l(\mathbf{x};\boldsymbol{\theta})\}
  -\max_{m\notin\{j,l\}}q_m(\mathbf{x};\boldsymbol{\theta}),
 \label{eq:activity-margin}
\end{equation}
where the maximum over an empty set is $-\infty$.

\begin{assumption}[Local stochastic regularity]
\label{ass:local-stochastic}
There exist a neighbourhood $U$ of $\boldsymbol{\theta}^*$, constants
$C,t_0>0$, and a measurable envelope $L$ with
$P_0L^{2+\eta}<\infty$ for some $\eta>0$ such
that:
\begin{enumerate}
 \item for all $\boldsymbol{\theta},\boldsymbol{\vartheta}\in U$ and every $j$,
 \[
  |q_j(\mathbf{X};\boldsymbol{\theta})-q_j(\mathbf{X};\boldsymbol{\vartheta})|
  +\|\nabla_{\boldsymbol{\theta}}q_j(\mathbf{X};\boldsymbol{\theta})-\nabla_{\boldsymbol{\theta}}q_j(\mathbf{X};\boldsymbol{\vartheta})\|
  \leq L(\mathbf{X})\|\boldsymbol{\theta}-\boldsymbol{\vartheta}\|
 \]
 almost surely, with the analogous local bound for $\nabla_{\boldsymbol{\theta}}^2q_j$;
 \item for every $j<l$ and $0<t<t_0$,
 \[
  P_0\{|g_{jl}(\mathbf{X};\boldsymbol{\theta}^*)|\leq tL(\mathbf{X}),\
          \Delta_{jl}(\mathbf{X};\boldsymbol{\theta}^*)\geq-tL(\mathbf{X})\}\leq Ct;
 \]
 \item simultaneous $t$-neighbourhoods of two distinct active contrasts have
 probability $O(t^2)$; and
 \item $\mathbf{A}=-\nabla_{\boldsymbol{\theta}}^2M(\boldsymbol{\theta}^*)$ is positive definite.
\end{enumerate}
\end{assumption}

The second condition is a local margin condition.  It follows under the
regular-face assumptions when the distribution of an active contrast has a
bounded density at zero, with suitable integrability for $L$.  The third
condition rules out first-order probability mass at multiple-face junctions.

\begin{lemma}[Quadratic mean expansion of the max criterion]
\label{lem:l2-expansion}
Under Assumptions~\ref{ass:boundary-geometry} and
\ref{ass:local-stochastic}, as $\mathbf{u}\to0$,
\begin{equation}
 \left\|m_{\boldsymbol{\theta}^*+\mathbf{u}}-m_{\boldsymbol{\theta}^*}-\mathbf{u}^\mathsf T\boldsymbol{\psi}^*\right\|_{P_0,2}
 =O(\|\mathbf{u}\|^{3/2})+O(\|\mathbf{u}\|^2).
 \label{eq:l2-expansion}
\end{equation}
The bound holds uniformly over $\mathbf{u}/\|\mathbf{u}\|$.
\end{lemma}

\begin{proof}
Put $r_{\mathbf u}=m_{\boldsymbol\theta^*+\mathbf u}-
m_{\boldsymbol\theta^*}-\mathbf u^\mathsf T\boldsymbol\psi^*$ and let
$T_{\mathbf u}$ be the union of the active tubes in
Assumption~\ref{ass:local-stochastic}(2), with $t=C_0\|\mathbf u\|$.
On $T_{\mathbf u}^c$ the winner is unchanged, and Taylor's theorem gives
$|r_{\mathbf u}|\leq C L\|\mathbf u\|^2$. On $T_{\mathbf u}$, the Lipschitz
bound gives $|r_{\mathbf u}|\leq C L\|\mathbf u\|$. After the standard
truncation implied by $P_0L^{2+\eta}<\infty$, the margin condition yields
\begin{align*}
 P_0r_{\mathbf u}^2
 &\leq C\|\mathbf u\|^4P_0L^2
   +C\|\mathbf u\|^2P_0(L^2\ind\{T_{\mathbf u}\})\\
 &=O(\|\mathbf u\|^4)+O(\|\mathbf u\|^3).
\end{align*}
Taking square roots proves \eqref{eq:l2-expansion}. Multiple-tube intersections
are $O(\|\mathbf u\|^2)$ by part (3), so they do not alter the bound. A full
bracketing construction and the precise truncation argument appear in
Supplementary Lemma 1.3.
\end{proof}

\begin{theorem}[Local empirical-process expansion]
\label{thm:local-empirical-expansion}
Under Assumptions~\ref{ass:boundary-geometry} and
\ref{ass:local-stochastic}, for every finite $R$,
\begin{multline}
 \sup_{\|\mathbf{u}\|\leq R}\left|
 n\left[(M_n-M)(\boldsymbol{\theta}^*+\mathbf{u}/\sqrt n)
       -(M_n-M)(\boldsymbol{\theta}^*)\right]
 -\mathbf{u}^\mathsf T\mathbb G_n\boldsymbol{\psi}^*
 \right|\\
 \xrightarrow{p}0.
 \label{eq:local-empirical-expansion}
\end{multline}
Consequently,
\begin{equation}
 n\{M_n(\boldsymbol{\theta}^*+\mathbf{u}/\sqrt n)-M_n(\boldsymbol{\theta}^*)\}
 =\mathbf{u}^\mathsf T\mathbb G_n\boldsymbol{\psi}^*
  -\tfrac12\mathbf{u}^\mathsf TAu+o_p(1),
 \label{eq:local-asymptotic-quadratic}
\end{equation}
uniformly on compact subsets of $\mathbb R^d$.
\end{theorem}

\begin{proof}
Write $r_{n,\mathbf u}=r_{\mathbf u/\sqrt n}$. Directly,
\begin{align*}
 &n(\mathbb P_n-P_0)\{m_{\boldsymbol\theta^*+\mathbf u/\sqrt n}
                       -m_{\boldsymbol\theta^*}\}
 -\mathbf u^\mathsf T\mathbb G_n\boldsymbol\psi^*\\
 &\hspace{35mm}=\sqrt n\,\mathbb G_n r_{n,\mathbf u}.
\end{align*}
Lemma~\ref{lem:l2-expansion} gives
$\sup_{\|\mathbf u\|\leq R}\|r_{n,\mathbf u}\|_{P_0,2}=O(n^{-3/4})$.
Supplementary Lemma 1.3 bounds the bracketing entropy by
$C d\log(C/\epsilon)$; the bracketing maximal inequality therefore gives
$\sup_{\|\mathbf u\|\leq R}|\sqrt n\mathbb G_nr_{n,\mathbf u}|=o_p(1)$.
This proves \eqref{eq:local-empirical-expansion}. Finally, Taylor expansion of
$M$ gives, uniformly on compact sets,
\[
 n\{M(\boldsymbol\theta^*+\mathbf u/\sqrt n)-M(\boldsymbol\theta^*)\}
 =\tfrac12\mathbf u^\mathsf T\mathbf H\mathbf u+o(1)
 =-\tfrac12\mathbf u^\mathsf T\mathbf A\mathbf u+o(1),
\]
which yields \eqref{eq:local-asymptotic-quadratic}
\citep{vandervaart1998,van1996weak}.
\end{proof}

\subsection{Rate, asymptotic linearity, and normality}

Define
\begin{equation}
 \mathbf{J}=P_0\{\boldsymbol{\psi}^*(\mathbf{X})\boldsymbol{\psi}^*(\mathbf{X})^\mathsf T\}.
 \label{eq:score-variance}
\end{equation}
At the interior maximizer, Theorem~\ref{thm:population-curvature} gives
$P_0\boldsymbol{\psi}^*=0$.

\begin{theorem}[Root-$n$ rate]
\label{thm:root-n-rate}
Suppose the conditions of Theorems~\ref{thm:cml-consistency} and
\ref{thm:local-empirical-expansion} hold.  If the approximate-maximization
error in \eqref{eq:approximate-maximizer} satisfies $nr_n=o_p(1)$, then
\begin{equation}
 \sqrt n(\widehat{\boldsymbol{\theta}}_n-\boldsymbol{\theta}^*)=O_p(1).
 \label{eq:root-n-rate}
\end{equation}
\end{theorem}

\begin{proof}
Let $\lambda_0=\lambda_{\min}(\mathbf A)>0$. For sufficiently small
$\delta$, the population expansion implies
$M(\boldsymbol\theta^*+\mathbf v)-M(\boldsymbol\theta^*)
\leq-(\lambda_0/4)\|\mathbf v\|^2$ whenever $\|\mathbf v\|\leq\delta$.
On the dyadic shell
$2^jC<\|\mathbf u\|\leq2^{j+1}C$, the local empirical term is
$O_p(2^jC)$ whereas the negative drift is at most
$-(\lambda_0/4)2^{2j}C^2$. Hence, uniformly over shells inside the
$\delta$-neighbourhood,
\[
 P\!\left(\sup_{\|\mathbf u\|>C}
 n\{M_n(\boldsymbol\theta^*+\mathbf u/\sqrt n)-M_n(\boldsymbol\theta^*)\}
 \geq-nr_n\right)\longrightarrow0
\]
as first $n\to\infty$ and then $C\to\infty$. Consistency excludes the
complement of that neighbourhood, and $nr_n=o_p(1)$ does not change the shell
comparison. Thus $\|\sqrt n(\widehat{\boldsymbol\theta}_n-
\boldsymbol\theta^*)\|=O_p(1)$.
\end{proof}

\begin{theorem}[Asymptotic linearity and normality]
\label{thm:asymptotic-normality}
Under the conditions of Theorem~\ref{thm:root-n-rate},
\begin{equation}
 \sqrt n(\widehat{\boldsymbol{\theta}}_n-\boldsymbol{\theta}^*)
 =\mathbf{A}^{-1}\mathbb G_n\boldsymbol{\psi}^*+o_p(1).
 \label{eq:asymptotic-linearity}
\end{equation}
If $\mathbf{J}$ is finite, then
\begin{equation}
 \sqrt n(\widehat{\boldsymbol{\theta}}_n-\boldsymbol{\theta}^*)
 \rightsquigarrow N_d(0,\mathbf{V}),
 \qquad \mathbf{V}=\mathbf{A}^{-1}\mathbf{J}\mathbf{A}^{-\mathsf T}.
 \label{eq:asymptotic-normality}
\end{equation}
This result remains centered on the CML functional $\boldsymbol{\theta}^*$ when the
ordinary mixture model is misspecified.
\end{theorem}

\begin{proof}
By Theorem~\ref{thm:root-n-rate}, the local parameter
$\widehat{\mathbf{u}}_n=\sqrt n(\widehat{\boldsymbol{\theta}}_n-\boldsymbol{\theta}^*)$ is tight. Put
$\mathbf Z_n=\mathbb G_n\boldsymbol\psi^*$. Completing the square in the local
criterion gives
\[
 \mathbf u^\mathsf T\mathbf Z_n-	frac12\mathbf u^\mathsf T\mathbf A\mathbf u
 =\tfrac12\mathbf Z_n^\mathsf T\mathbf A^{-1}\mathbf Z_n
 -\tfrac12(\mathbf u-\mathbf A^{-1}\mathbf Z_n)^\mathsf T
 \mathbf A(\mathbf u-\mathbf A^{-1}\mathbf Z_n).
\]
Uniformity on compact sets, tightness, and $nr_n=o_p(1)$ force
$\widehat{\mathbf u}_n-\mathbf A^{-1}\mathbf Z_n=o_p(1)$, proving
\eqref{eq:asymptotic-linearity}. Since
$\mathbf Z_n\rightsquigarrow N_d(\mathbf0,\mathbf J)$ by the multivariate
central limit theorem, Slutsky's theorem gives covariance
$\mathbf A^{-1}\mathbf J\mathbf A^{-\mathsf T}$.
\end{proof}

\begin{proposition}[Transfer to a local algorithmic solution]
\label{prop:algorithmic-transfer}
Let $\widetilde{\boldsymbol{\theta}}_n\to_p\boldsymbol{\theta}^*$ be an interior solution produced by an
algorithm.  Suppose ties occur with probability zero and
\begin{equation}
 \mathbb P_n\boldsymbol{\psi}(\mathbf{X};\widetilde{\boldsymbol{\theta}}_n)=o_p(n^{-1/2}),
 \label{eq:algorithmic-score-error}
\end{equation}
and suppose the local score process is stochastically differentiable at
$\boldsymbol{\theta}^*$ with derivative $-\mathbf{A}$.  Then
\begin{equation}
 \sqrt n(\widetilde{\boldsymbol{\theta}}_n-\boldsymbol{\theta}^*)
 =\mathbf{A}^{-1}\mathbb G_n\boldsymbol{\psi}^*+o_p(1).
 \label{eq:algorithmic-linearity}
\end{equation}
Thus a consistent stationary solution has the same first-order distribution
as the local CML maximizer when its optimization error is smaller than the
statistical score error.
\end{proposition}

\begin{proof}
Stochastic differentiability gives
\[
 \mathbb P_n\boldsymbol{\psi}(\mathbf{X};\widetilde{\boldsymbol{\theta}}_n)
 =\mathbb P_n\boldsymbol{\psi}^*-\mathbf{A}(\widetilde{\boldsymbol{\theta}}_n-\boldsymbol{\theta}^*)
  +o_p(n^{-1/2}+\|\widetilde{\boldsymbol{\theta}}_n-\boldsymbol{\theta}^*\|).
\]
Use $P_0\boldsymbol{\psi}^*=0$, multiply by $\sqrt n$, and solve the resulting linear
equation.  Positive definiteness of $\mathbf{A}$ and consistency absorb the remainder.
\end{proof}

\begin{theorem}[Validity of the nonparametric bootstrap]
\label{thm:bootstrap-validity}
Let $\mathbb P_n^{\#}$ be the empirical measure of a size-$n$ sample drawn
with replacement from $\mathbb P_n$, and let
$\widehat{\boldsymbol{\theta}}_n^{\#}$ be a bootstrap approximate maximizer
initialized in the same consistent basin as
$\widehat{\boldsymbol{\theta}}_n$.  Suppose, conditionally on the data, its
optimization error is $o_p(n^{-1})$, and suppose the localized classes in
Theorem~\ref{thm:local-empirical-expansion} are bootstrap-Donsker and
$\mathbf{J}$ is positive definite.  Then
\begin{equation}
 \sup_{C\in\mathcal C}
 \left|
 P^{\#}\!\left\{
 \sqrt n(\widehat{\boldsymbol{\theta}}_n^{\#}
          -\widehat{\boldsymbol{\theta}}_n)\in C
 \right\}
 -P\{N_d(\mathbf{0},\mathbf{V})\in C\}
 \right|\xrightarrow{p}0,
 \label{eq:bootstrap-validity}
\end{equation}
where $\mathcal C$ is the collection of convex Borel subsets of
$\mathbb R^d$.  Thus the ordinary nonparametric bootstrap consistently
estimates the first-order distribution, provided each bootstrap fit reaches
the relevant local solution.
\end{theorem}

\begin{proof}
Let $\mathbb G_n^\#=\sqrt n(\mathbb P_n^\#-\mathbb P_n)$. Bootstrap
equicontinuity of the localized remainder class gives, conditionally in
probability and uniformly for bounded $\mathbf u$,
\begin{align*}
 &n\{M_n^\#(\widehat{\boldsymbol\theta}_n+\mathbf u/\sqrt n)
       -M_n^\#(\widehat{\boldsymbol\theta}_n)\}\\
 &=\mathbf u^\mathsf T\mathbb G_n^\#
   \boldsymbol\psi(\cdot;\widehat{\boldsymbol\theta}_n)
   -\tfrac12\mathbf u^\mathsf T\widehat{\mathbf A}_n\mathbf u
   +o_{p^\#}(1).
\end{align*}
Here $\widehat{\mathbf A}_n\to_p\mathbf A$, and the conditional central limit
theorem gives
$\mathbb G_n^\#\boldsymbol\psi(\cdot;\widehat{\boldsymbol\theta}_n)
\rightsquigarrow_{p}N_d(\mathbf0,\mathbf J)$. Completing the square, as in the
proof of Theorem~\ref{thm:asymptotic-normality}, and using the conditional
$o_p(n^{-1})$ optimization error yields
\[
 \sqrt n(\widehat{\boldsymbol\theta}_n^\#-\widehat{\boldsymbol\theta}_n)
 =\widehat{\mathbf A}_n^{-1}\mathbb G_n^\#
 \boldsymbol\psi(\cdot;\widehat{\boldsymbol\theta}_n)+o_{p^\#}(1).
\]
The conditional Slutsky theorem gives the Gaussian limit with covariance
$\mathbf V$. Since that law is nondegenerate, its boundary assigns zero mass
to convex-set continuity boundaries, yielding the displayed uniform metric
\citep{van1996weak}.
\end{proof}


\subsection{Feasible boundary-curvature estimation}
\label{sec:estimation}
The surface representation in \eqref{eq:boundary-curvature} is useful for
theory but direct numerical integration becomes impractical as the observation
dimension $p$ grows.  We instead estimate its equivalent Dirac-delta
representation using observations from $P_0$.  This avoids both a spatial grid
and the incorrect substitution of the fitted CML mixture for the unknown
data-generating density.

For $j<l$, define the active-pair indicator
\begin{equation}
 \chi_{jl}(\mathbf{x};\boldsymbol{\theta})
 =\ind\!\left\{
 \min(q_j(\mathbf{x};\boldsymbol{\theta}),q_l(\mathbf{x};\boldsymbol{\theta}))
 >\max_{m\notin\{j,l\}}q_m(\mathbf{x};\boldsymbol{\theta})
 \right\},
 \label{eq:active-pair-indicator}
\end{equation}
with the indicator equal to one when $k=2$.  Let $K$ be a kernel and $h=h_n$
a bandwidth.  The proposed estimator is
\begin{equation}
 \widehat{\mathbf{B}}_n
 =\frac{1}{nh}
 \sum_{i=1}^n\sum_{1\leq j<l\leq k}
 K\!\left\{\frac{g_{jl}(\mathbf{X}_i;\widehat{\boldsymbol{\theta}}_n)}{h}\right\}
 \chi_{jl}(\mathbf{X}_i;\widehat{\boldsymbol{\theta}}_n)
 \nabla_{\boldsymbol{\theta}}g_{jl}(\mathbf{X}_i;\widehat{\boldsymbol{\theta}}_n)
 \nabla_{\boldsymbol{\theta}}g_{jl}(\mathbf{X}_i;\widehat{\boldsymbol{\theta}}_n)^\mathsf T.
 \label{eq:boundary-estimator}
\end{equation}
Its cost is linear in $n$ and quadratic in $k$, and it does not require a grid
in $\mathbb R^p$.

\begin{assumption}[Kernel and tube regularity]
\label{ass:kernel-tube}
\begin{enumerate}
 \item $K$ is bounded, Lipschitz, symmetric, compactly supported, and
 $\int K(v)\,dv=1$.
 \item $h\to0$, $nh\to\infty$, and
 $\|\widehat{\boldsymbol{\theta}}_n-\boldsymbol{\theta}^*\|/h=o_p(1)$.
 \item In neighbourhoods of active faces, the conditional surface moments of
 $\nabla_{\boldsymbol{\theta}}g_{jl}\nabla_{\boldsymbol{\theta}}g_{jl}^\mathsf T$ given $g_{jl}=v$ are continuous at
 $v=0$ and dominated by an integrable envelope.
 \item Contributions from $h$-neighbourhoods of multiple-tie junctions are
 $o(1)$ after the $h^{-1}$ normalization.
\end{enumerate}
\end{assumption}

For a root-$n$ consistent estimator, the plug-in condition in part (2) follows
from $\sqrt n\,h\to\infty$.  The familiar choice $h\asymp n^{-1/5}$ satisfies
the three rate restrictions, although data-driven bandwidth selection and
sensitivity analysis remain necessary in applications.

To describe precision, define the oracle estimator
$\widetilde{\mathbf{B}}_n$ by replacing
$\widehat{\boldsymbol{\theta}}_n$ with $\boldsymbol{\theta}^*$ in
\eqref{eq:boundary-estimator}.  For an active pair, let
\begin{equation}
 \boldsymbol{\Gamma}_{jl}(v)
 =\int_{\{\mathbf{x}:g_{jl}(\mathbf{x};\boldsymbol{\theta}^*)=v\}}
 \frac{p_0(\mathbf{x})\chi_{jl}(\mathbf{x};\boldsymbol{\theta}^*)
 \mathbf{W}_{jl}(\mathbf{x})}
 {\|\nabla_{\mathbf{x}}g_{jl}(\mathbf{x};\boldsymbol{\theta}^*)\|}
 \,d\mathcal H^{p-1}(\mathbf{x}),
 \label{eq:surface-moment}
\end{equation}
where
\begin{equation}
 \mathbf{W}_{jl}(\mathbf{x})
 =\nabla_{\boldsymbol{\theta}}g_{jl}(\mathbf{x};\boldsymbol{\theta}^*)
  \nabla_{\boldsymbol{\theta}}g_{jl}(\mathbf{x};\boldsymbol{\theta}^*)^\mathsf T.
 \label{eq:boundary-outer-product}
\end{equation}

\begin{theorem}[Bias, variance, and oracle limit distribution]
\label{thm:boundary-oracle-clt}
In addition to Assumption~\ref{ass:kernel-tube}, suppose each
$\boldsymbol{\Gamma}_{jl}$ is twice continuously differentiable at zero,
$K$ is a second-order kernel, and the corresponding fourth-order surface
moments are finite.  Let
$\mu_2(K)=\int v^2K(v)\,dv$ and $R(K)=\int K(v)^2\,dv$.  Then
\begin{align}
 E\widetilde{\mathbf{B}}_n
 &=\mathbf{B}+\frac{h^2\mu_2(K)}2
   \sum_{j<l}\boldsymbol{\Gamma}_{jl}''(0)+o(h^2),
 \label{eq:boundary-bias}\\
 \operatorname{Var}\{\operatorname{vec}(\widetilde{\mathbf{B}}_n)\}
 &=\frac1{nh}\boldsymbol{\Omega}_{\mathbf{B}}
   +o\{(nh)^{-1}\},
 \label{eq:boundary-variance}
\end{align}
where
\begin{equation}
\begin{aligned}
 \boldsymbol{\Omega}_{\mathbf{B}}
 &=R(K)\sum_{j<l}\int_{\mathcal F_{jl}^*}
 \frac{p_0(\mathbf{x})}
 {\|\nabla_{\mathbf{x}}g_{jl}(\mathbf{x};\boldsymbol{\theta}^*)\|}\\
 &\quad\times
 \operatorname{vec}\{\mathbf{W}_{jl}(\mathbf{x})\}
 \operatorname{vec}\{\mathbf{W}_{jl}(\mathbf{x})\}^\mathsf T
 \,d\mathcal H^{p-1}(\mathbf{x}).
\end{aligned}
 \label{eq:boundary-asymptotic-variance}
\end{equation}
If $nh^5\to0$, then
\begin{equation}
 \sqrt{nh}\,\operatorname{vec}
 (\widetilde{\mathbf{B}}_n-\mathbf{B})
 \rightsquigarrow N_{d^2}(\mathbf{0},\boldsymbol{\Omega}_{\mathbf{B}}).
 \label{eq:boundary-oracle-clt}
\end{equation}
Without undersmoothing, the same conclusion holds after centering by the
leading bias in \eqref{eq:boundary-bias}.
\end{theorem}

\begin{proof}
The coarea formula writes the expectation as
$\sum_{j<l}\int K(t)\boldsymbol{\Gamma}_{jl}(ht)\,dt$.
A second-order Taylor expansion gives \eqref{eq:boundary-bias}.  Applying the
same representation to the squared vectorized summand gives
\eqref{eq:boundary-variance}; cross-pair terms are lower order because
first-order mass at multiple-face junctions is excluded.  A triangular-array
Lindeberg argument yields the centered normal limit.  Finally,
$\sqrt{nh}h^2\to0$ is equivalent to $nh^5\to0$ and removes the leading bias.
\end{proof}

\subsection{Finite-sample curvature stabilization}

Kernel variability can make
$\widehat{\mathbf{A}}_n=\widehat{\mathbf{I}}_{c,n}-
\widehat{\mathbf B}_n$ nonpositive in finite samples even when
$\mathbf{A}$ is positive definite.  To make this failure visible while still
providing usable inference, define
\begin{align}
 \widehat\lambda_n
 &=\lambda_{\max}\!\left(
 \widehat{\mathbf I}_{c,n}^{-1/2}
 \widehat{\mathbf B}_n
 \widehat{\mathbf I}_{c,n}^{-1/2}\right),                 \label{eq:relative-boundary-eigenvalue}\\
 \widehat\gamma_n
 &=\min\left\{1,\frac{1-\varepsilon_n}{\widehat\lambda_n}\right\},
 \qquad
 \widehat{\mathbf B}_n^{\mathrm{stab}}
 =\widehat\gamma_n\widehat{\mathbf B}_n,                 \label{eq:boundary-stabilization}
\end{align}
where $\varepsilon_n\downarrow0$.  We use
$\varepsilon_n=n^{-1/4}$ in the simulations.

\begin{proposition}[Asymptotic inactivity of stabilization]
\label{prop:stabilization-inactive}
Suppose $\mathbf{I}_c$ and
$\mathbf{A}=\mathbf{I}_c-\mathbf{B}$ are positive definite and the conditions of
Theorem~\ref{thm:sandwich-consistency} hold.  If
$\varepsilon_n\to0$, then
\begin{equation}
 P(\widehat\gamma_n=1)\to1,
 \qquad
 \widehat{\mathbf{B}}_n^{\mathrm{stab}}\to_p\mathbf{B}.
 \label{eq:stabilization-inactive}
\end{equation}
Consequently, replacing $\widehat{\mathbf B}_n$ by its stabilized version does
not change the first-order limiting distribution.
\end{proposition}

\begin{proof}
Positive definiteness of $\mathbf{A}$ implies
\[
 \lambda_{\max}(\mathbf{I}_c^{-1/2}\mathbf{B}
                         \mathbf{I}_c^{-1/2})<1.
\]
Consistency of $\widehat{\mathbf I}_{c,n}$ and
$\widehat{\mathbf B}_n$, together with continuity of inverse square roots and
eigenvalues on the positive-definite cone, gives
$\widehat\lambda_n\to_p\lambda_{\max}<1$.  Since
$\varepsilon_n\to0$, the inequality
$\widehat\lambda_n\leq1-\varepsilon_n$ holds with probability tending to one.
Hence $\widehat\gamma_n=1$ with probability tending to one, proving both
claims.
\end{proof}

\begin{proposition}[Plug-in error rate]
\label{prop:boundary-plugin-rate}
Under the conditions of Theorem~\ref{thm:boundary-oracle-clt}, suppose the
probability contribution of $h$-neighbourhoods of multiple-tie junctions,
after $h^{-1}$ normalization, is $O(\zeta_h)$ with $\zeta_h\to0$.  Then
\begin{equation}
 \|\widehat{\mathbf{B}}_n-\mathbf{B}\|_F
 =O_p\!\left(
 h^2+(nh)^{-1/2}
 +\frac{\|\widehat{\boldsymbol{\theta}}_n
             -\boldsymbol{\theta}^*\|}{h}
 +\zeta_h
 \right).
 \label{eq:boundary-plugin-rate}
\end{equation}
For root-$n$ consistent $\widehat{\boldsymbol{\theta}}_n$ and
$h\asymp n^{-1/5}$, all displayed terms converge to zero.
\end{proposition}

\begin{proof}
The first two terms follow from
Theorem~\ref{thm:boundary-oracle-clt}.  Lipschitz continuity of $K$, the
$h^{-1}$ scaling in its argument, and local derivative envelopes bound the
parameter-substitution error by the third term.  Classification of an active
pair can differ only near a multiple-tie junction or within a tube whose width
is proportional to the parameter error; these contributions are controlled by
$\zeta_h$ and the same substitution term.
\end{proof}

\begin{theorem}[Consistency of the boundary-curvature estimator]
\label{thm:boundary-estimator-consistency}
Under Assumptions~\ref{ass:boundary-geometry},
\ref{ass:local-stochastic}, and \ref{ass:kernel-tube},
\begin{equation}
 \widehat{\mathbf{B}}_n\xrightarrow{p}\mathbf{B}.
 \label{eq:boundary-estimator-consistency}
\end{equation}
\end{theorem}

\begin{proof}
First fix $\boldsymbol{\theta}=\boldsymbol{\theta}^*$.  The coarea formula gives, for each active pair,
\begin{multline*}
 \frac1hP_0\left[
 K\{g_{jl}(\mathbf{X};\boldsymbol{\theta}^*)/h\}\chi_{jl}(\mathbf{X};\boldsymbol{\theta}^*)
 \nabla_{\boldsymbol{\theta}}g_{jl}(\mathbf{X};\boldsymbol{\theta}^*)\nabla_{\boldsymbol{\theta}}g_{jl}(\mathbf{X};\boldsymbol{\theta}^*)^\mathsf T
 \right]\\
 =\int K(v/h)h^{-1}
 \int_{\{g_{jl}=v\}}
 \frac{p_0(\mathbf{x})\chi_{jl}(\mathbf{x};\boldsymbol{\theta}^*)
 \nabla_{\boldsymbol{\theta}}g_{jl}(\mathbf{x};\boldsymbol{\theta}^*)\nabla_{\boldsymbol{\theta}}g_{jl}(\mathbf{x};\boldsymbol{\theta}^*)^\mathsf T}
 {\|\nabla_{\mathbf{x}}g_{jl}(\mathbf{x};\boldsymbol{\theta}^*)\|}
 \,d\mathcal H^{p-1}(\mathbf{x})\,dv.
\end{multline*}
Continuity of the inner surface moment and $\int K=1$ make this converge to
the $(j,l)$ term of \eqref{eq:boundary-curvature}.  The variance is
$O\{(nh)^{-1}\}$ under the envelope conditions, so the centered empirical
term vanishes.  The condition
$\|\widehat{\boldsymbol{\theta}}_n-\boldsymbol{\theta}^*\|/h=o_p(1)$, Lipschitz regularity, and negligible
multiple-tie neighbourhoods make substitution of $\widehat{\boldsymbol{\theta}}_n$ and the
estimated active indicator asymptotically negligible.  Summing over the
finitely many pairs proves the result.
\end{proof}

The choice $h\asymp n^{-1/5}$ is used here for consistency and conventional
MSE-type bias--variance behaviour, not for an undersmoothed centred CLT.

Define the feasible fixed-classification and score matrices
\begin{align}
 \widehat{\mathbf{I}}_{c,n}
 &=-\frac1n\sum_{i=1}^n\sum_{j=1}^k
 \ind\{a_{\widehat{\boldsymbol{\theta}}_n}(\mathbf{X}_i)=j\}
 \nabla_{\boldsymbol{\theta}}^2q_j(\mathbf{X}_i;\widehat{\boldsymbol{\theta}}_n),                         \label{eq:ic-estimator}\\
 \widehat{\mathbf{J}}_n
 &=\frac1n\sum_{i=1}^n
 \widehat{\boldsymbol{\psi}}_i\widehat{\boldsymbol{\psi}}_i^\mathsf T,
 \qquad
 \widehat{\boldsymbol{\psi}}_i=\boldsymbol{\psi}(\mathbf{X}_i;\widehat{\boldsymbol{\theta}}_n),              \label{eq:j-estimator}\\
 \widehat{\mathbf{A}}_n
 &=\widehat{\mathbf{I}}_{c,n}-\widehat{\mathbf{B}}_n,               \label{eq:a-estimator}\\
 \widehat{\mathbf{V}}_n
 &=\widehat{\mathbf{A}}_n^{-1}\widehat{\mathbf{J}}_n\widehat{\mathbf{A}}_n^{-\mathsf T}.
 \label{eq:v-estimator}
\end{align}

\begin{theorem}[Consistency of feasible sandwich inference]
\label{thm:sandwich-consistency}
Suppose the assumptions of
Theorem~\ref{thm:boundary-estimator-consistency} hold, the interior Hessian and
score classes satisfy uniform laws of large numbers near $\boldsymbol{\theta}^*$, and $\mathbf{A}$
is nonsingular.  Then
\begin{equation}
 \widehat{\mathbf{I}}_{c,n}\to_p \mathbf{I}_c,\qquad
 \widehat{\mathbf{J}}_n\to_p \mathbf{J},\qquad
 \widehat{\mathbf{A}}_n\to_p \mathbf{A},\qquad
 \widehat{\mathbf{V}}_n\to_p \mathbf{V}.
 \label{eq:sandwich-consistency}
\end{equation}
Consequently, $\widehat{\mathbf{V}}_n/n$ consistently estimates the first-order
covariance matrix of $\widehat{\boldsymbol{\theta}}_n$.
\end{theorem}

\begin{proof}
Consistency of $\widehat{\boldsymbol{\theta}}_n$, the local uniform laws, and the fact that
ties have probability zero give consistency of $\widehat{\mathbf{I}}_{c,n}$ and
$\widehat{\mathbf{J}}_n$.  Theorem~\ref{thm:boundary-estimator-consistency} gives the
remaining component of $\widehat{\mathbf{A}}_n$.  Matrix inversion is continuous on the
set of nonsingular matrices, so the final assertion follows from the
continuous mapping theorem.
\end{proof}

\begin{corollary}[First-order equivalence after stabilization]
\label{cor:stabilized-equivalence}
Under Proposition~\ref{prop:stabilization-inactive} and
Theorem~\ref{thm:sandwich-consistency}, let
$\widehat{\mathbf V}_n^{\mathrm{stab}}$ be obtained from
$\widehat{\mathbf B}_n^{\mathrm{stab}}$.  Then
$P(\widehat{\mathbf V}_n^{\mathrm{stab}}=\widehat{\mathbf V}_n)\to1$ and
$\widehat{\mathbf V}_n^{\mathrm{stab}}\to_p\mathbf V$.  Hence stabilized and
unstabilized Wald statistics have the same first-order limit whenever
stabilization is asymptotically inactive.
\end{corollary}

\begin{corollary}[Wald inference]
\label{cor:wald-inference}
Let $\mathbf{R}$ be a fixed $r\times d$ matrix of rank $r$ and suppose
$\mathbf{R}\mathbf{V}\mathbf{R}^\mathsf T$ is nonsingular.  Under
$H_0:\mathbf{R}\boldsymbol{\theta}^*=\mathbf{c}$,
\begin{equation}
 n(\mathbf{R}\widehat{\boldsymbol{\theta}}_n-\mathbf{c})^\mathsf T
 (\mathbf{R}\widehat{\mathbf{V}}_n\mathbf{R}^\mathsf T)^{-1}
 (\mathbf{R}\widehat{\boldsymbol{\theta}}_n-\mathbf{c})
 \rightsquigarrow\chi_r^2.
 \label{eq:wald-statistic}
\end{equation}
In particular, componentwise intervals based on
$\{\widehat{\mathbf{V}}_{n,ss}/n\}^{1/2}$ have asymptotic coverage equal to their
nominal level.
\end{corollary}

\begin{corollary}[Smooth transformations]
\label{cor:delta-method}
Let $\boldsymbol{\tau}:\mathbb R^d\to\mathbb R^s$ be continuously
differentiable at $\boldsymbol{\theta}^*$ with Jacobian matrix
$\mathbf{D}_{\tau}$.  Then
\begin{equation}
 \sqrt n\{\boldsymbol{\tau}(\widehat{\boldsymbol{\theta}}_n)
 -\boldsymbol{\tau}(\boldsymbol{\theta}^*)\}
 \rightsquigarrow
 N_s(\mathbf{0},\mathbf{D}_{\tau}\mathbf{V}
                    \mathbf{D}_{\tau}^\mathsf T),
 \label{eq:delta-method}
\end{equation}
and the plug-in covariance obtained by replacing
$\mathbf{D}_{\tau}$ and $\mathbf{V}$ by consistent estimators is consistent.
This result covers scientifically interpretable transformations of mixture
weights, locations, covariance parameters, and component contrasts.
\end{corollary}


\subsection{Computation and implementation}
\label{sec:computation}
The implementation uses one distribution adapter for every component family.
An adapter defines the unconstrained coordinates, inverse transformation, and
log density used by fitting, classification, differentiation, and inference.
Consequently, the estimated scores and decision contrasts cannot silently use
a density different from the fitted model.  Positive-definite covariance
matrices use log-Cholesky coordinates, and mixture weights use
reference-category multinomial logits.

\begin{center}
\fbox{\begin{minipage}{0.94\linewidth}
\textbf{Algorithm 1: Feasible boundary-aware covariance estimation}
\begin{enumerate}
 \item Fit the hard-classification likelihood to obtain
 $\widehat{\boldsymbol\theta}_n$.
 \item Compute the fitted labels.
 \item Compute $\widehat{\mathbf I}_{c,n}$.
 \item Compute $\widehat{\mathbf J}_n$.
 \item Compute every pairwise contrast $g_{jl}$ and its parameter gradient.
 \item Estimate $\widehat{\mathbf B}_n$ by \eqref{eq:boundary-estimator}.
 \item Form $\widehat{\mathbf A}_n=\widehat{\mathbf I}_{c,n}-
 \widehat{\mathbf B}_n$ and stabilize it if necessary, recording the diagnostic.
 \item Report $\widehat{\mathbf V}_n=\widehat{\mathbf A}_n^{-1}
 \widehat{\mathbf J}_n\widehat{\mathbf A}_n^{-\mathsf T}$ and standard errors
 from $\widehat{\mathbf V}_n/n$.
\end{enumerate}
\end{minipage}}
\end{center}

For the multivariate skew-normal component, we use
\begin{equation}
 f(\mathbf{x};\boldsymbol{\xi},\boldsymbol{\Omega},\boldsymbol{\alpha})
 =2\phi_p(\mathbf{x};\boldsymbol{\xi},\boldsymbol{\Omega})
 \Phi\!\left\{
 \boldsymbol{\alpha}^{\mathsf T}\boldsymbol{\Omega}^{-1/2}
 (\mathbf{x}-\boldsymbol{\xi})
 \right\},
 \label{eq:skew-normal-density}
\end{equation}
where $\boldsymbol{\Omega}^{-1/2}$ is the symmetric inverse square root.  The
same expression is used by the fitting library and the inference code.

Hard-EM fitting explicitly requests the hard mode.  We use both K-means and
Gaussian-mixture initializations.  When component families are distinct, all
family-to-component assignments are fitted and the result with the largest
empirical classification likelihood is retained.  The objective gap, selected
initialization, component weights, tail parameter, covariance eigenvalues, and
skewness norm are retained as basin-stability diagnostics.  The result is then
mapped to the manuscript's canonical family order.  Repeated families are
ordered by a prespecified fitted location functional.  Neither operation uses
the generating parameters or simulated labels.

The reference implementation evaluates derivatives by centered finite
differences in the common unconstrained coordinate system.  Evaluations are
vectorized over observations, so a full contrast-gradient matrix requires
two score evaluations per parameter coordinate rather than two evaluations
per observation and coordinate.  A deliberately slow observation-wise
implementation is retained as an independent numerical check.

The contrast $g_{jl}$ is a dimensionless log posterior odds.  Our prespecified
pilot bandwidth therefore uses its natural scale,
\begin{equation}
 h_{jl}=1.06n^{-1/5}.
 \label{eq:pilot-bandwidth}
\end{equation}
The global empirical spread of $g_{jl}$ is reported but does not multiply the
bandwidth: in separated mixtures it is dominated by observations far from the
boundary and produces severe oversmoothing. This is a pilot rule, not an
optimality claim. Every analysis reports $h_{jl}$, the robust global spread,
the number of observations for which $(j,l)$ is active, and the number falling
inside the kernel support.
Pairs with insufficient effective observations are flagged rather than
silently assigned a wider bandwidth.  The simulation study examines
multiplier sensitivity around \eqref{eq:pilot-bandwidth}.

We report the raw corrected matrix and its nonpositive-curvature failure rate.
Primary finite-sample inference uses the eigenvalue-constrained version in
\eqref{eq:boundary-stabilization}, with $\varepsilon_n=n^{-1/4}$.  The
associated shrinkage factor is reported for every replication; the correction
is asymptotically unmodified by Proposition~\ref{prop:stabilization-inactive}.

Population CML targets used in simulation are approximated from several
independent large reference samples.  We retain every resulting target vector
and report the between-reference Monte Carlo standard error of each coordinate.
This prevents numerical target uncertainty from being treated as zero.  Target
fits and ordinary simulation fits use the same multistart hard-EM,
classification-objective selection, and canonical-ordering rules.

The nonparametric bootstrap used to assess local inferential accuracy is
initialized at the original estimate and constrained to a local trust region.
Replicates reaching that region's boundary are reported as failures.  Fresh
global bootstrap fits that reselect the family assignment answer a different
question---stability under model and basin selection---and are reported
separately rather than pooled into a Gaussian standard error.


\section{Results}
\label{sec:results}

\subsection{Simulation study}
\label{sec:simulations}
We used two-component Normal--Student-$t$ and Normal--skew-normal mixtures in
dimensions $p=2$ and $p=5$, with Euclidean location separations representing
separated, moderate-overlap, and strong-overlap regimes.  Balanced designs used
$n=500$ and $n=2000$; additional experiments considered weights $(0.2,0.8)$,
three heterogeneous components, and $p=10$.  There were 1000 replications per
primary cell and 500 in the dimension stress cells, for 31,000 attempted fits.
Bandwidth multipliers $0.75$, $1$, and $1.5$ were evaluated without refitting.
All random seeds, including failed replications, were retained.

Figure~\ref{fig:simulation-profiles} shows the principal two-dimensional
designs before aggregation over regimes.  The correction has little effect in
the separated cells, but the gap between the methods grows as the active
boundary moves into a region of greater probability mass.

\begin{figure*}[t]
\centering
\includegraphics[width=\textwidth]{figures/simulation_profiles.pdf}
\caption{Standard-error calibration and 95\% coverage in the principal
two-dimensional simulations. Shapes distinguish separation regimes; solid
blue and dashed grey curves denote stabilized boundary-aware and naive
inference. Each point is a coordinate median over usable replications.}
\label{fig:simulation-profiles}
\end{figure*}

The estimand was the population classification maximum-likelihood functional,
not the generating ordinary-mixture parameter.  For each of 18 scenarios we
used 10 independent reference samples of size 100,000, reduced to five for the
initial $p=10$ stress calculation.  All reference fits succeeded.  Across the
18 targets the largest coordinate Monte Carlo standard error was 0.033 and the
largest range across independent reference fits was 0.322.  No selected target
fit violated the prespecified admissibility diagnostics.

Table~\ref{tab:simulation-main} reports coordinate-median coverage and the
coordinate-median ratio of estimated standard error to empirical standard
deviation.  The stabilized method uses the central bandwidth multiplier.  The
correction is most consequential under overlap and imbalance.  For example, in
the Normal--Student-$t$, $p=2$, moderate-overlap design with $n=500$, median coverage
increased from 0.886 to 0.925 and the median standard-error ratio from 0.758 to
0.920.  Under imbalance the corresponding coverage increased from 0.834 to
0.895.  Under separation and $n=2000$, both methods are closer, as expected.

\begin{table}[t]
\centering
\caption{Selected simulation results. Entries are coordinate-median estimated
standard error divided by empirical standard deviation, followed in
parentheses by coordinate-median 95\% coverage.}
\label{tab:simulation-main}
\begin{tabular}{llrrcc}
\toprule
Families & Regime & $p$ & $n$ & Naive & Stabilized \\
\midrule
Normal--Student-$t$ & moderate & 2 & 500  & 0.758 (0.886) & 0.920 (0.925) \\
Normal--Student-$t$ & moderate & 2 & 2000 & 0.827 (0.896) & 0.994 (0.940) \\
Normal--Student-$t$ & imbalanced & 2 & 500 & 0.615 (0.834) & 0.844 (0.895) \\
Normal--Student-$t$ & strong & 2 & 500 & 0.491 (0.750) & 0.675 (0.844) \\
Normal--Student-$t$ & separated & 5 & 2000 & 0.956 (0.939) & 1.018 (0.950) \\
Normal--skew-normal & strong & 2 & 2000 & 0.573 (0.868) & 0.762 (0.922) \\
Normal--skew-normal & moderate & 5 & 2000 & 0.892 (0.933) & 0.994 (0.943) \\
\bottomrule
\end{tabular}
\end{table}

Table~\ref{tab:simulation-quantiles} aggregates all retained coordinate-level
primary results rather than selecting favourable coordinates.  Coordinate-level
results, including the difficult coordinates, are reported in the supplement;
no difficult coordinate was removed from the analysis.

\begin{table}[t]
\centering
\caption{Across-coordinate distribution of calibration results in the primary
study.  Entries are the 25th percentile, median, and 75th percentile across all
scenario--sample-size--coordinate cells; corrected results use the central
bandwidth.}
\label{tab:simulation-quantiles}
\begin{tabular}{lrrrrrr}
\toprule
& \multicolumn{3}{c}{SE/empirical SD} & \multicolumn{3}{c}{Empirical coverage}\\
Method & Q25 & Median & Q75 & Q25 & Median & Q75\\
\midrule
Naive & 0.654 & 0.867 & 0.952 & 0.878 & 0.917 & 0.935\\
Corrected, stabilized & 0.780 & 0.991 & 1.032 & 0.911 & 0.937 & 0.947\\
\bottomrule
\end{tabular}
\end{table}

Figure~\ref{fig:simulation-calibration} summarizes every primary scenario.
Nearly all points lie above the equality line: the correction increases both
the estimated-standard-error ratio and coverage, generally toward the ideal
horizontal reference rather than indiscriminately beyond it.

\begin{figure*}[t]
\centering
\includegraphics[width=\textwidth]{figures/simulation_calibration.pdf}
\caption{Scenario-level comparison of naive and stabilized boundary-aware
inference across the complete primary study. Point colour denotes the family
design and point size distinguishes larger sample sizes. The dashed diagonal
denotes equality; dotted lines mark ideal SE calibration and nominal coverage.}
\label{fig:simulation-calibration}
\end{figure*}

Across the primary study, 4214 scenario--replication cases had at least one raw
nonpositive curvature result among the three prespecified bandwidths.  This is
why stabilization is reported as a diagnostic rather than hidden as a purely
numerical adjustment.  Raw correction frequently produced a nonpositive estimated curvature, with the
frequency increasing under strong overlap and in higher dimension.  In
contrast, the stabilized correction was computable whenever the underlying
fit and naive information calculation succeeded.  Stabilization was active in
10.5\% of Normal--Student-$t$, $p=2$, moderate-overlap replications at $n=500$, falling
to 3.4\% at $n=2000$; in the separated $p=2$ cells it fell from 5.1\% to 1.0\%.
Under the regular theory, stabilization frequency should decrease with sample
size; these paired simulation cells show precisely that tendency.

Normal--Student-$t$ fitting succeeded in every primary replication.  Operational fit
failure for models involving a skew-normal component ranged from zero to 7.6\%,
with the maximum occurring in the $p=10$ stress design.  We did not replace or
discard these seeds.  Conditional summaries use all replications for which the
base fit and naive information were available, and failure rates are reported
separately.

A post-confirmatory convergence extension used 500 replications at $n=8000$
for the two $p=2$ strong-overlap models and at $n=4000$ for the two
three-component models.  Stabilized median coverage reached 0.940 for the
Normal--skew-normal strong-overlap case and 0.942 for the $p=5$ three-component
case.  The Normal--Student-$t$ strong-overlap case improved to 0.923, compared with
0.721 for naive inference.  The $p=2$ three-component case remained difficult,
with stabilized median coverage 0.863.  This last result illustrates slow
finite-sample convergence when several decision faces and family-assignment
uncertainty interact; it is reported as a limitation rather than removed from
the design.

Figure~\ref{fig:extension-boundary} places this extension beside a separate
boundary-estimator convergence experiment. The latter shows decreasing RMSE
from sample size 250 to 16,000, directly checking the feasible surface-moment
estimator independently of the Wald calculations.
The supplement additionally uses all replication-level estimator vectors to
check standardized root-$n$ quantiles, Kolmogorov distance from Gaussianity,
scaled RMSE, and scaled bias. Thus consistency, rate, asymptotic normality,
variance calibration, coverage, feasible-curvature convergence, stabilization,
and bootstrap behaviour each receive a distinct empirical diagnostic.

\begin{figure*}[t]
\centering
\includegraphics[width=\textwidth]{figures/extension_boundary_convergence.pdf}
\caption{Large-sample and estimator-convergence diagnostics. Left:
coordinate-median coverage in four post-confirmatory cells; segments join
naive and corrected results. Right: RMSE of a scalar boundary-curvature
functional, with an $n^{-1/2}$ reference slope.}
\label{fig:extension-boundary}
\end{figure*}


\subsection{Real data analysis}
\label{sec:applications}
We applied the method to the Dry Bean data of \citet{koklu2020drybean},
which contain 12 measured dimensions and four shape descriptors extracted from
images of 13,611 individual grains belonging to seven registered varieties.
The class labels were not used in estimation.  They were retained only to give
an external assessment of the fitted partition.  We standardized all 16
features and used the first two principal-component scores, thereby avoiding a
choice of variables based on the reported class labels.

Before running the bootstrap, we screened six data configurations (breast
cancer, wine, handwritten digits, and three pairs of bean varieties), two
dimensions, and Normal--Student-$t$ and Normal--skew-normal specifications.
The screen used numerical admissibility, component weights, effective kernel
support near the estimated boundary, and whether the boundary correction was
material.  It was used to select one regular illustration and one deliberately
difficult diagnostic case, rather than to select the largest apparent
improvement in clustering accuracy.  The complete screening results and failed
fits are retained in the reproduction archive.

For each selected data set we fitted a two-component Normal--Student-$t$
classification mixture.  We compared the analytic standard errors with 500
nonparametric bootstrap samples optimized locally within a coordinate-wise
trust region around the original solution.  This is the bootstrap analogue of
the local target in our theory.  We also performed 200 unrestricted refits from
the full fitting procedure.  The latter diagnostic allows the selected basin
to change and hence measures instability that is outside the local Gaussian
approximation.  Table~\ref{tab:applications} summarizes the results.

Figure~\ref{fig:application-boundaries} displays the standardized scores,
coloured by external variety labels, together with the fitted classification
boundary. Labels played no role in preprocessing or fitting.

\begin{figure*}[t]
\centering
\includegraphics[width=\textwidth]{figures/application_boundaries.pdf}
\caption{Dry Bean principal-component scores and fitted Normal--Student-$t$
classification boundaries. Colours show external labels used only for
evaluation; black curves are equality sets of the fitted weighted log densities.}
\label{fig:application-boundaries}
\end{figure*}

\begin{table}[t]
\centering
\scriptsize
\setlength{\tabcolsep}{3pt}
\caption{Dry Bean applications. The two final columns give the coordinate-
median ratio of an analytic standard error to the standard deviation from
successful local bootstrap fits. Boundary support is the number of observations
inside the fitted kernel tube. Local success is the percentage of 500 locally
optimized bootstrap refits that passed all convergence diagnostics.}
\label{tab:applications}
\resizebox{\linewidth}{!}{%
\begin{tabular}{lrrrrrr}
\toprule
Varieties & Sample size & Adjusted Rand & Kernel count & Local success & Naive ratio & Corrected ratio \\
\midrule
Dermason--Seker & 5573 & 0.867 & 33 & 95.0\% & 0.798 & 0.991 \\
Cali--Barbunya  & 2952 & 0.282 & 44 & 52.4\% & 0.538 & 0.850 \\
\bottomrule
\end{tabular}
}

\smallskip
\begin{minipage}{0.94\linewidth}\small
The median stabilized-to-naive standard-error inflation was 1.284 and 1.650,
respectively.  The percentages of unrestricted refits remaining within a
maximum coordinate distance of 0.75 from the original solution were 59.0\%
and 46.5\%, respectively.
\end{minipage}
\end{table}

\paragraph{A regular validation case.}
For Dermason--Seker the fitted partition agrees closely with the independently
recorded varieties (adjusted Rand index 0.867), both estimated component
weights are substantial, and 475 of 500 local bootstrap fits succeeded.  The
naive standard errors are systematically too small: their median ratio to the
local bootstrap standard deviation is 0.798.  Boundary-aware standard errors
raise this ratio to 0.991.  At the coordinate level the stabilized ratios range
from 0.886 to 1.082, whereas the naive ratios range from 0.576 to 0.939.  Thus
the aggregate improvement is not driven by one parameter.  The estimated
boundary curvature required no finite-sample shrinkage in this fit.

\begin{table}[t]
\centering
\scriptsize
\setlength{\tabcolsep}{3pt}
\caption{Dermason--Seker inference for interpretable fitted quantities. The
location distance is measured in standardized principal-component space.
Here ``s.e.'' means standard error and ``bootstrap s.d.'' means the standard
deviation of successful local-bootstrap estimates.}
\label{tab:dermason-interpretable}
\resizebox{\linewidth}{!}{%
\begin{tabular}{lrrrrl}
\toprule
Quantity & Estimate & Naive s.e. & Corrected s.e. & Bootstrap s.d. & Corrected 95\% interval\\
\midrule
Mixing weight $\pi_1$ & 0.345 & 0.006 & 0.008 & 0.008 & (0.329, 0.361)\\
$\|\boldsymbol\mu_1-\boldsymbol\mu_2\|$ & 5.202 & 0.036 & 0.037 & 0.035 & (5.130, 5.273)\\
Student-$t$ degrees of freedom $\nu$ & 20.680 & 3.848 & 4.184 & 5.533 & (12.479, 28.881)\\
\bottomrule
\end{tabular}
}
\end{table}

Scientifically, the correction changes uncertainty most clearly for the fitted
prevalence of the first hard-classification group: its standard error increases
by about 31\%, bringing it close to the local-bootstrap scale.  The estimated
location separation is so pronounced that its uncertainty changes little, as
predicted by the low-boundary-density argument.  For the Student-$t$ tail
parameter, the corrected interval is wider than the naive interval, although
the remaining gap to the local-bootstrap standard deviation warns that tail
estimation is harder.  The correction therefore changes substantive precision
selectively; it does not inflate every uncertainty by a common factor.

Figure~\ref{fig:application-calibration} confirms that the aggregate result is
not an artefact of taking medians. Dermason--Seker corrected ratios cluster
around one across the coordinate vector; Cali--Barbunya shows a substantial
but incomplete improvement for most coordinates.

\begin{figure*}[t]
\centering
\includegraphics[width=\textwidth]{figures/application_se_calibration.pdf}
\caption{Coordinate-level analytic standard errors divided by local-bootstrap
standard deviations. The dotted reference at one represents exact local
calibration.}
\label{fig:application-calibration}
\end{figure*}

Only 59.0\% of unrestricted bootstrap refits remained in the original local
basin.  Location-coordinate standard deviations from unrestricted refitting
were consequently orders of magnitude larger for some coordinates.  This does
not contradict the local calibration result: it identifies a separate layer
of optimization and selection uncertainty.  In practice, the local standard
errors should therefore be reported together with the basin-stability
diagnostic.

The unrestricted distances in Figure~\ref{fig:application-basins} make this
distinction visible. Both examples have a local group of refits, but
Dermason--Seker also has a separate group near distance five and
Cali--Barbunya has a broad nonlocal tail. A global standard deviation would
conflate these discrete changes with regular local sampling variation.

\begin{figure*}[t]
\centering
\includegraphics[width=\textwidth]{figures/application_basin_diagnostics.pdf}
\caption{Maximum coordinate distance between each unrestricted bootstrap
refit and the original solution. The dashed line is the prespecified local
trust-region radius 0.75.}
\label{fig:application-basins}
\end{figure*}

\paragraph{A nonregular diagnostic case.}
The Cali--Barbunya pair has stronger overlap, as reflected in the lower
external adjusted Rand index (0.282) and the larger median analytic correction
(a factor of 1.650).  The stabilized standard errors again move substantially
toward the local-bootstrap scale, improving the median ratio from 0.538 to
0.850.  However, only 262 of 500 local optimizations succeeded, and only 46.5\%
of unrestricted refits stayed near the original solution.  These diagnostics
fail the operational stability requirement needed to treat the local normal
approximation as a complete account of uncertainty.  We therefore present
this example as a stress test, not as evidence for routine Wald inference.
Together, the two applications show both the practical gain from estimating
boundary curvature and the need to distinguish local sampling uncertainty
from basin-selection instability.


\section{Discussion}
\label{sec:discussion}
Hard classification creates a distinctive inferential problem: fitted labels
can be locally constant for almost every individual observation even though
their population partition changes at first order.  The missing curvature is
concentrated on decision hypersurfaces.  Our surface representation makes this
mechanism explicit, and its positive semidefinite sign explains why treating
the labels as fixed tends to produce standard errors that are too small.

The feasible estimator turns that geometric identity into a practical
procedure.  It smooths only a scalar log-score contrast, regardless of the
dimension of the observations, and it samples the unknown data distribution
directly.  This is preferable to integrating a fitted mixture density over an
estimated surface, which would impose correct specification precisely where
the theory does not require it.  The simulations and the Dermason--Seker
analysis show that the correction can improve local calibration without the
cost of a large bootstrap.  Under clear separation the correction contracts,
as predicted by the surface formula.

\paragraph{What the method does not do.}
The method does not target the ordinary mixture parameter: its estimand is the
population hard-classification optimum.  It does not solve uncertainty from
label selection or component-family selection.  Nor does it guarantee valid
inference near nonregular configurations with multiple active boundaries.
Regular active faces,
a unique local target, and positive population curvature are substantive
conditions.  Nearly tangent boundaries, appreciable multiple ties, vanishing
components, or weakly identified skewness can invalidate the approximation or
slow its convergence.  Global family assignment and optimization-basin
selection are discrete sources of uncertainty.  The unrestricted bootstrap
results, particularly for Cali--Barbunya, show that they should not be hidden
inside a Wald standard error.
The method estimates local sampling uncertainty conditional on the selected
basin; it does not account for global optimization-basin selection.  When
unrestricted bootstrap refits often move to another basin, corrected Wald
intervals should be reported only as conditional local summaries.

These qualifications suggest a reporting protocol.  Alongside corrected
standard errors, an analysis should provide component weights, boundary kernel
support, the raw relative boundary eigenvalue, any stabilization factor,
multistart objective gaps, and the fraction of local bootstrap fits that pass
their convergence diagnostics.  When unrestricted refits frequently change
basin, conditional local intervals remain interpretable for the selected
solution but should not be presented as a complete uncertainty statement.

Several extensions are promising.  Uniform inference as two population maxima
approach one another would require a nonregular, set-valued limit theory.
Joint selection and estimation could instead use selective or subsampling
methods.  Higher-order or locally adaptive bandwidths may improve estimation
when boundary density varies sharply, while automatic differentiation could
replace finite differences in large parameter spaces.  More broadly, the same
surface-curvature principle applies to other estimators built from maxima of
smooth criteria.  Developing this connection beyond mixture classification is
an important direction for future work.

\section{Conclusion}
The inferential consequence of hard classification is geometric: decision
boundaries move with the fitted parameter and contribute first-order
curvature even though the boundaries themselves have probability zero. The
proposed surface correction estimates this missing curvature with a
one-dimensional kernel calculation and yields materially better calibrated
local uncertainty where fixed-label inference is most anti-conservative. The
theory, simulations, and applications also delimit the claim: corrected Wald
inference describes a regular local optimum and cannot absorb instability from
competing optimization basins or nearly nonregular classification faces.
Reporting the correction together with boundary support, stabilization, and
basin diagnostics therefore supplies both an inferential method and an
explicit warning when a Gaussian local summary is insufficient.

\section*{Data and code availability}
The Dry Bean data are publicly available from the UCI Machine Learning
Repository under a Creative Commons Attribution 4.0 licence
\citep{uci2020drybean}. Code, prespecified screening output, raw job manifests,
and scripts producing every table and figure are maintained at
\url{https://github.com/samyajoypal/CEM_Asymptotics}.

\section*{Conflict of interest}
The authors declare that they have no conflict of interest.

\section*{Funding}
No external funding was received for this study.

\section*{Author contributions}
S.P. contributed conceptualization, methodology, software, formal analysis,
visualization, and preparation of the original draft. R.E.H. contributed
review and editing of the manuscript. Both authors approved the final
manuscript.
\bibliographystyle{abbrvnat}
\bibliography{references}
\end{document}
